
\documentclass{article}
\usepackage{colortbl}
\usepackage{makecell}
\usepackage{multirow}
\usepackage{supertabular}

\begin{document}

\newcounter{utterance}

\centering \large Interaction Transcript for game `codenames', experiment `abstract', episode 0 with reflexion{-}llama{-}3.1{-}8b{-}instruct.
\vspace{24pt}

{ \footnotesize  \setcounter{utterance}{1}
\setlength{\tabcolsep}{0pt}
\begin{supertabular}{c@{$\;$}|p{.15\linewidth}@{}p{.15\linewidth}p{.15\linewidth}p{.15\linewidth}p{.15\linewidth}p{.15\linewidth}}
    \# & $\;$A & \multicolumn{4}{c}{Game Master} & $\;\:$B\\
    \hline

    \theutterance \stepcounter{utterance}  
    & & & \multicolumn{2}{p{0.3\linewidth}}{
        \cellcolor[rgb]{0.9,0.9,0.9}{
            \makecell[{{p{\linewidth}}}]{
                \texttt{\tiny{[GM$|$GM]}}
                {'hidden': {'team': ['affirmation', 'lunacy', 'adventurousness', 'interconnection', 'oppressiveness', 'perpetration', 'anticommunism'], 'innocent': ['fantasy', 'endearment', 'manner'], 'opponent': ['unification', 'lessening'], 'assassin': ['optimism']}, 'revealed': {'team': {'team': [], 'innocent': [], 'opponent': [], 'assassin': []}, 'opponent': {'team': [], 'innocent': [], 'opponent': [], 'assassin': []}}}
            }
        }
    }
    & & \\ \\

    \theutterance \stepcounter{utterance}  
    & & \multicolumn{4}{p{0.6\linewidth}}{
        \cellcolor[rgb]{0.9,0.9,0.9}{
            \makecell[{{p{\linewidth}}}]{
                \texttt{\tiny{[P1$\langle$GM]}}
                \texttt{You have attempted this same task before and did not succeed. These are the plans you wrote after those attempts.} \\
\texttt{Trial \#1: New plan: } \\
\\ 
\texttt{In my last attempt, I provided a clue that was related to three words, but I didn't check if the clue was semantically related to all of them. In this case, "Extremity" is more closely related to "lunacy" and "grandiosity" than to "oppressiveness". This led to the Game Master rejecting my message because it didn't follow the specified format. Next time, I will make sure to provide a clue that is semantically related to all the target words, and I will also double{-}check the format of my response. For example, I could provide a clue like "CLUE: Chaos" and "TARGETS: lunacy, oppressiveness, grandiosity" to relate all three words to a common theme.} \\
\\ 
\texttt{Trial \#2: New plan: } \\
\\ 
\texttt{In my last attempt, I provided a clue that was related to three words, but I didn't check if the clue was semantically related to all of them. In this case, "Extremity" is more closely related to "lunacy" and "grandiosity" than to "oppressiveness". This led to the Game Master rejecting my message because it didn't follow the specified format. Next time, I will make sure to provide a clue that is semantically related to all the target words, and I will also double{-}check the format of my response. } \\
\\ 
\texttt{To achieve this, I will carefully analyze the relationships between the words and choose a clue that is common to all of them. I will also make sure to check the list of remaining words for my opponent and the distractor words to avoid providing a clue that is too obvious or too vague. For example, I could provide a clue like "CLUE: Chaos" and "TARGETS: lunacy, oppressiveness, grandiosity" to relate all three words to a common theme.} \\
\\ 
\\ 
\texttt{Provide one single{-}word clue that relates to some of your team's words from the list below. You can choose to how many and to which words your clue relates to.} \\
\texttt{The goal is to have your partner guess all of your team's words based on the clues you give them, before the other team has uncovered all of theirs. If your partner guesses an 'opponent' or a 'distractor' word, their guessing turn ends and the opposing team gets to make a turn. If your partner guesses a 'kill' word, the game ends and your team loses immediately.} \\
\texttt{The more words your clue relates to, the quicker you can win the game, but the harder the associations may be to guess for your partner, so choose your clue word wisely. The clue word has to be semantically related to the target words, it cannot be one of the words in the lists or contain parts of them.} \\
\\ 
\texttt{Always give your single{-}word clue and your comma{-}separated list of related target words in the following format and make your answers as short as possible, never include any other text than is required in this form:} \\
\\ 
\texttt{CLUE: <WORD>} \\
\texttt{TARGETS: <TARGETS>} \\
\\ 
\texttt{Your team words are: affirmation, lunacy, adventurousness, interconnection, oppressiveness, perpetration, anticommunism.} \\
\texttt{Your opponent's team words are: unification, lessening.} \\
\texttt{Distractor words are: fantasy, endearment, manner.} \\
\texttt{The kill words are: optimism.} \\
            }
        }
    }
    & \\ \\

    \theutterance \stepcounter{utterance}  
    & \multicolumn{4}{p{0.6\linewidth}}{
        \cellcolor[rgb]{0.9,0.9,0.9}{
            \makecell[{{p{\linewidth}}}]{
                \texttt{\tiny{[P1$\rangle$GM]}}
                \texttt{CLUE: Excess} \\
\texttt{TARGETS: lunacy, oppressiveness, grandiosity} \\
            }
        }
    }
    & & \\ \\

    \theutterance \stepcounter{utterance}  
    & & & \multicolumn{2}{p{0.3\linewidth}}{
        \cellcolor[rgb]{0.9,0.9,0.9}{
            \makecell[{{p{\linewidth}}}]{
                \texttt{\tiny{[GM$|$GM]}}
                {'player': 'Cluegiver', 'type': 'target is hallucination', 'utterance': 'CLUE: Excess\nTARGETS: lunacy, oppressiveness, grandiosity', 'message': "Targeted word 'grandiosity' was not listed, you can only target words provided in the lists.", 'target': 'grandiosity', 'previous_targets': [], 'remaining_words': ['unification', 'optimism', 'endearment', 'perpetration', 'affirmation', 'adventurousness', 'manner', 'lessening', 'fantasy', 'lunacy', 'interconnection', 'oppressiveness', 'anticommunism']}
            }
        }
    }
    & & \\ \\

    \theutterance \stepcounter{utterance}  
    & \multicolumn{4}{p{0.6\linewidth}}{
        \cellcolor[rgb]{0.95,0.95,0.95}{
            \makecell[{{p{\linewidth}}}]{
                \texttt{\tiny{[P2$\rightarrow$P1]}}
                \texttt{Post{-}episode Reflexion (not sent to GM):} \\
\texttt{New plan: } \\
\\ 
\texttt{In my last attempts, I provided clues that were related to three words, but I didn't check if the clue was semantically related to all of them. This led to the Game Master rejecting my message because it didn't follow the specified format. Next time, I will make sure to provide a clue that is semantically related to all the target words, and I will also double{-}check the format of my response. } \\
\\ 
\texttt{To achieve this, I will carefully analyze the relationships between the words and choose a clue that is common to all of them. I will also make sure to check the list of remaining words for my opponent and the distractor words to avoid providing a clue that is too obvious or too vague. } \\
\\ 
\texttt{I will start by identifying the common theme or characteristic among the target words, such as "excess" or "intensity". Then, I will choose a clue that is closely related to this theme and is not too obvious or too vague. For example, I could provide a clue like "CLUE: Excess" and "TARGETS: lunacy, oppressiveness, grandiosity" to relate all three words to a common theme.} \\
            }
        }
    }
    & & \\ \\

\end{supertabular}
}

\end{document}
